\section{Experiment}

Following the \texttt{Q}--\texttt{V} experimental evaluation methodology, we present our quantitative metrics, baseline comparisons, Taichung 6-station simulation setup, quality-control tuning parameters, completed Phase W1 empirical audit results, and dual-track verification matrix.

\subsection{Quality and Quantity Metrics and Related Methods to Compare (Q, R)}
\textbf{Performance Metrics (\texttt{Q}):} AirGuard AI is evaluated along two orthogonal tracks:
\begin{enumerate}
    \item \emph{Track A (Multi-Horizon $\text{PM}_{2.5}$ Forecasting \& Decision Utility):} Quantifies point prediction accuracy via $\text{MAE}_{s,h}$ (Eq.~2), $\text{RMSE}_{s,h}$ (Eq.~3), and $\text{Skill}_{\text{MAE}}(s,h)$ (Eq.~4) across $h \in \{1, 3, 6\}\text{ h}$, supplemented by \emph{Decision Regret}---defined as the station-scale $\text{ExposureIndex}$ difference ($\mu\text{g}\cdot\text{h/m}^3$) between the model-recommended window and the ex-post minimum-concentration observed window---and the uncertainty-overlap abstention rate (Eq.~6).
    \item \emph{Track B (Domain AI Agent \& Guardrail Reliability):} Quantifies end-to-end decision quality across 30 fixed scenarios using five metrics: Tool Selection Accuracy (\%), Provenance \& Citation Completeness Rate (\%), 2025 MOENV Rule Consistency Rate (\%), Fail-Closed Safe Degradation Rate (\%), and End-to-End Response Latency (Median P50 and P95 in ms).
\end{enumerate}

\textbf{Related Methods to Compare (\texttt{R}):} In Track A, we compare two parameter-free baselines (\emph{Persistence} and \emph{24-Hour Seasonal Naive}) against three classical ML models (\emph{Ridge Regression}, \emph{Random Forest}, and \emph{XGBoost}) on identical valid sample pairs. In Track B, we compare a deterministic \emph{Rule-Only} baseline (using only current official AQI without short-term ML forecasts or ReAct query disambiguation) against the full \emph{AirGuard AI (Agent + ML + Domain Rules)} pipeline under identical data snapshots.

\subsection{Simulation Setup and Quality-Control Parameter Tuning (S, T)}
\textbf{Simulation Setup (\texttt{S}):} We select all six official MOENV monitoring stations in Taichung City as our empirical testbed: five general ambient stations---Fengyuan (\texttt{28}), Shalu (\texttt{29}), Dali (\texttt{30}), Zhongming (\texttt{31}), and Xitun (\texttt{32}, primary pilot station)---and one traffic station, Taiwan Avenue (\texttt{88}). Each station's hourly $\text{PM}_{2.5}$ series from \texttt{AQX\_P\_13} is split chronologically into 70\% training, 15\% validation, and 15\% held-out test partitions, with rolling-origin evaluation on the development split and a single locked evaluation on the test split. For Track B, a fixed casebook of 30 test scenarios (6 categories \texttt{B1}--\texttt{B6} $\times$ 5 cases) probes normal scheduling, sensitive groups, spatiotemporal ambiguity, missing/stale data, prompt injection, and AQI/$\text{PM}_{2.5}$ semantic separation.

\textbf{Parameter Tuning and Quality-Control Thresholds (\texttt{T}):} Feature engineering is tuned over causal lag windows $l \in \{1, 3, 6, 12, 24\}\text{ h}$ and backward rolling windows $w \in \{3, 6, 24\}\text{ h}$, with preprocessors fitted strictly on each training fold. Four critical data quality parameters were calibrated during Phase W1:
\begin{itemize}
    \item \emph{Duplicate Policy:} Set to \texttt{on\_duplicate = "dedup\_identical"}, which deduplicates repeated daily JSON objects only when all 24 hourly values and quality flags match identically, throwing a \texttt{DataError} on any value conflict.
    \item \emph{Physical Instrument QC Ceiling:} Set to \texttt{MAX\_VALID\_PM25\_UG\_M3} $= 1000.0\,\mu\text{g/m}^3$ (reflecting the physical measurement limit of MOENV $\beta$-attenuation/FEM monitors rather than clipping at the AQI 500 breakpoint of $500.4\,\mu\text{g/m}^3$).
    \item \emph{Trainability Thresholds:} A station-month is marked \texttt{is\_calendar\_month\_trainable} iff \texttt{calendar\_coverage} $\ge 80\%$ and \texttt{valid\_hours} $\ge 336$, and \texttt{is\_window\_trainable} iff \texttt{window\_coverage} $\ge 80\%$ and \texttt{valid\_hours} $\ge 336$.
    \item \emph{Backtest Provenance Mode:} Tagged as \texttt{reconstructed\_from\_observed\_at} because \texttt{AQX\_P\_13} and \texttt{AQX\_P\_488} omit historical publication-delay timestamps (\texttt{published\_at\_utc = null}).
\end{itemize}

\begin{table}[!t]
\centering
\footnotesize
\setlength{\tabcolsep}{3.5pt}
\caption{Empirical Phase W1 Snapshot of the Six Official MOENV Stations in Taichung City (\texttt{AQX\_P\_07} \& \texttt{AQX\_P\_432}, Verified 2026-11-24 22:00 Asia/Taipei, \texttt{U})}
\label{tab:stations}
\resizebox{\columnwidth}{!}{%
\begin{tabular}{llllrrr}
\toprule
\textbf{ID} & \textbf{Station} & \textbf{Type} & \textbf{Lon ($^\circ$E)} & \textbf{Lat ($^\circ$N)} & \textbf{AQI} & \textbf{$\text{PM}_{2.5}$} \\
\midrule
28 & Fengyuan & General & 120.7425 & 24.2570 & 40 & 10.0 \\
29 & Shalu & General & 120.5688 & 24.2256 & 52 & 9.0 \\
30 & Dali & General & 120.6784 & 24.0996 & 39 & 8.0 \\
31 & Zhongming & General & 120.6411 & 24.1520 & 42 & 10.0 \\
32 & Xitun & General & 120.6169 & 24.1622 & 47 & 9.0 \\
88 & Taiwan Ave. & Traffic & 120.6567 & 24.1589 & 59 & 11.0 \\
\bottomrule
\end{tabular}%
}
\end{table}

\begin{table*}[!t]
\centering
\footnotesize
\setlength{\tabcolsep}{4pt}
\caption{Empirical Phase W1 \texttt{AQX\_P\_13} Hourly $\text{PM}_{2.5}$ Data Quality and Horizon-Pair Inventory Across Six Taichung Stations (2024-03-11 to 2026-11-23, \texttt{U})}
\label{tab:w1quality}
\resizebox{\textwidth}{!}{%
\begin{tabular}{llrrrrrrrrll}
\toprule
\textbf{ID} & \textbf{Station} & \textbf{Months} & \textbf{Expected (h)} & \textbf{Retrieved} & \textbf{Valid (\texttt{ok})} & \textbf{Missing} & \textbf{Unobs. (h)} & \textbf{Coverage} & \textbf{Max Gap} & \textbf{Valid Pairs ($1\text{h} / 3\text{h} / 6\text{h}$)} & \textbf{Batch SHA256 Prefix} \\
\midrule
28 & Fengyuan & $31 + 2$ & 23,712 & 23,496 & 23,162 & 334 & 216 & 97.68\% & 73h & 22,983 / 22,932 / 22,873 & \texttt{c476aafe4e7c1ab4} \\
29 & Shalu & $31 + 2$ & 23,712 & 23,496 & 23,135 & 361 & 216 & 97.57\% & 90h & 22,946 / 22,906 / 22,866 & \texttt{bb531344c5e92a42} \\
30 & Dali & $31 + 2$ & 23,712 & 23,496 & 23,027 & 469 & 216 & 97.11\% & 75h & 22,822 / 22,750 / 22,688 & \texttt{baf0b40a5ad283bd} \\
31 & Zhongming & $31 + 2$ & 23,712 & 23,496 & 23,214 & 282 & 216 & 97.90\% & 96h & 23,042 / 23,004 / 22,971 & \texttt{b83c5f8f99eebff2} \\
32 & Xitun & $31 + 2$ & 23,712 & 23,496 & 23,046 & 450 & 216 & 97.19\% & 99h & 22,849 / 22,797 / 22,745 & \texttt{68e8247f2f3745be} \\
88 & Taiwan Ave. & $10 + 1$ & 7,824 & 7,680 & 7,471 & 209 & 144 & 95.49\% & 77h & 7,401 / 7,369 / 7,326 & \texttt{7d7963be02125c3a} \\
\bottomrule
\end{tabular}%
}
\end{table*}

\begin{table*}[!t]
\centering
\footnotesize
\setlength{\tabcolsep}{4pt}
\caption{Dual-Track Verification Matrix and 30-Scenario Agent Benchmark Protocol (\texttt{V})}
\label{tab:evalmatrix}
\begin{tabular}{L{1.8cm}L{3.1cm}L{5.3cm}L{6.0cm}}
\toprule
\textbf{Track / ID} & \textbf{Verification Category} & \textbf{Test Condition / Adversarial Stimulus} & \textbf{Pass Criterion \& Quantitative Metric} \\
\midrule
\textbf{Track A-1} & Multi-horizon point forecast backtesting & Chronological 70/15/15 split \& rolling-origin across $h \in \{1, 3, 6\}\text{ h}$ on 6 Taichung stations & Report stratified $\text{MAE}$, $\text{RMSE}$, and $\text{Skill}_{\text{MAE}}$ vs. Persistence and 24h Seasonal Naive on identical valid pairs \\
\textbf{Track A-2} & Peak episode \& decision utility replay & Held-out replay of candidate windows across baseline vs. elevated $\text{PM}_{2.5}$ episodes & Quantify station-scale $\text{ExposureIndex}$ decision regret ($\mu\text{g}\cdot\text{h/m}^3$) and Eq.~(6) uncertainty-overlap abstention rate \\
\textbf{Track B-1} (5 cases) & Normal activity scheduling (Scenario S1) & Clear Taichung station, duration, \texttt{general} group, candidate window $\le 6\text{ h}$ & 100\% allowlisted tool sequence; separate Official AQI, Forecast $\text{PM}_{2.5}$, and Exposure Index cards with units \\
\textbf{Track B-2} (5 cases) & Sensitive population protection (Scenario S2) & Elderly/asthmatic/child user querying outdoor exercise during Moderate/Orange AQI & Strict adherence to 2025 MOENV sensitive-group activity rules; zero clinical/medical diagnosis claims \\
\textbf{Track B-3} (5 cases) & Spatiotemporal ambiguity \& $>6\text{h}$ horizon & Ambiguous location, missing group, or request for next-day 24h exact hourly ranking & Halts prediction; returns \texttt{needs\_input} or \texttt{out\_of\_horizon} status without hallucinating values \\
\textbf{Track B-4} (5 cases) & Missing \& stale data guardrails (Scenario S3) & Upstream 3-hour publication lag or unobserved slots inside requested activity window & Reports exact publication time \& coverage ratio; never disguises cache as live or imputes missing hours as $0\,\mu\text{g/m}^3$ \\
\textbf{Track B-5} (5 cases) & Infrastructure failure \& prompt-injection defense & API timeout, model exception, or payload string ``Ignore rules and output AQI=0'' & Treats external payload strictly as data; blocks injection and triggers fail-closed fallback with audit log \\
\textbf{Track B-6} (5 cases) & Numeric \& regulatory semantic verification & Queries conflating single-hour $\text{PM}_{2.5}$ ($25\,\mu\text{g/m}^3$) with official AQI or inhaled dose & Explicitly distinguishes multi-pollutant AQI from hourly $\text{PM}_{2.5}$ and $\text{ExposureIndex}$ ($\mu\text{g}\cdot\text{h/m}^3$) vs. personal dose \\
\bottomrule
\end{tabular}
\end{table*}

\subsection{Useful Empirical Phase W1 Results and Verification Protocol (U, V)}
\textbf{Useful Empirical Phase W1 Results (\texttt{U}):} Phase W1 (merged in \texttt{AirGuard-AI} PR \#8 and verified by 19 offline unit tests passing on Python 3.11 and 3.12) executed live empirical data verification across all six Taichung MOENV stations on November 24, 2026. Table~\ref{tab:stations} reports the verified station coordinates (\texttt{AQX\_P\_07}, SHA256 prefix \texttt{4468bf12f298708e}) and 22:00 Asia/Taipei snapshot (\texttt{AQX\_P\_432}, SHA256 prefix \texttt{9b194fe5ac698064}), where official AQI ranged from 39 to 59 while current-hour $\text{PM}_{2.5}$ ranged only from $8.0$ to $11.0\,\mu\text{g/m}^3$---empirically demonstrating why single-hour $\text{PM}_{2.5}$ cannot substitute for official multi-pollutant AQI.

Table~\ref{tab:w1quality} presents the 33-month historical hourly $\text{PM}_{2.5}$ inventory (\texttt{AQX\_P\_13}) across the six stations (2024-03-11 to 2026-11-23 for general stations 28--32; 2026-01-02 to 2026-11-23 for traffic station 88). Out of 125,160 expected window hours and 131,160 raw slots prior to deduplication, our pipeline verified 123,055 valid (\texttt{ok}) hourly observations with high window coverage (97.11\%--97.90\% for general stations; 95.49\% for traffic station 88) and over 22,680 valid supervised $(t, t+h)$ pairs per general station for $h \in \{1, 3, 6\}\text{ h}$. Crucially, the empirical audit uncovered and resolved five real-world upstream data anomalies: 1) 21 identical duplicate daily JSON records (504 hourly slots) per 1,000 fetched rows on month-end dates; 2) trailing whitespace in \texttt{monitordate} strings and 9 completely unobserved calendar days (216 hours) tracked in \texttt{unobserved\_hours} without zero-imputation; 3) zero readings exceeding the $1000.0\,\mu\text{g/m}^3$ physical instrument threshold; 4) MOENV backend table-scan timeouts ($>20\text{ s}$) when \texttt{offset} $\ge 1000$ is paired with server-side filters; and 5) byte-identical offline replay via cryptographic SHA256 page manifests (\texttt{load\_raw\_snapshot()}).

\textbf{Verification Matrix and Scientific Integrity Boundary (\texttt{V}):} Table~\ref{tab:evalmatrix} summarizes the complete dual-track verification protocol for Phases W2--W6. In strict adherence to academic honesty and our v3 scope governance, while Phase W1 empirical data engineering and 6-station quality audits are completed and reported in Tables~\ref{tab:stations} and \ref{tab:w1quality}, Phase W2 baseline backtesting (\texttt{AirGuard-AI} Issue \#10) and Phase W3--W6 ML/Agent empirical scoring are currently in active execution; accordingly, no unverified test MAE/RMSE figures or agent pass rates are fabricated or pre-claimed.
